\documentclass[12pt,a4paper]{article}

% ============================================
% PACKAGES
% ============================================
\usepackage[utf8]{inputenc}          % UTF-8 encoding
\usepackage[T1]{fontenc}             % Font encoding
\usepackage{lmodern}                 % Modern Latin fonts
\usepackage[english]{babel}          % Language settings
\usepackage{amsmath,amssymb,amsthm}  % Math packages
\usepackage{graphicx}                % Include images
\usepackage{hyperref}                % Clickable links
\usepackage{xcolor}                  % Colors
\usepackage{listings}                % Code listings
\usepackage{geometry}                % Page margins
\usepackage{fancyhdr}                % Headers and footers
\usepackage{lipsum}                  % Dummy text for demonstration

% ============================================
% DOCUMENT SETTINGS
% ============================================
\geometry{margin=1in}

% Hyperlink settings
\hypersetup{
    colorlinks=true,
    linkcolor=blue,
    filecolor=magenta,
    urlcolor=cyan,
    pdftitle={Sample LaTeX Document},
    pdfauthor={Your Name},
}

% Code listing settings
\lstset{
    basicstyle=\ttfamily\small,
    breaklines=true,
    frame=single,
    numbers=left,
    numberstyle=\tiny\color{gray},
    keywordstyle=\color{blue},
    commentstyle=\color{green!60!black},
    stringstyle=\color{red},
}

% Header/Footer
\pagestyle{fancy}
\fancyhf{}
\fancyhead[L]{Sample Document}
\fancyhead[R]{\today}
\fancyfoot[C]{\thepage}

% ============================================
% TITLE
% ============================================
\title{Sample \LaTeX{} Document}
\author{Your Name}
\date{\today}

% ============================================
% DOCUMENT BODY
% ============================================
\begin{document}

\maketitle

\begin{abstract}
This is a sample \LaTeX{} document demonstrating various features including mathematical equations, figures, tables, code listings, and bibliography references. Use this as a template to get started with your own documents. Now I have added more text to the abstract to ensure it spans multiple lines and provides a better overview of the document's content and purpose. Change 3. Change 4. Change 5.
\end{abstract}

\tableofcontents
\newpage

% --------------------------------------------
\section{Introduction}
% --------------------------------------------

Welcome to \LaTeX{}! This document serves as a starting point for your \LaTeX{} journey in VS Code. \LaTeX{} is a powerful typesetting system widely used in academia and scientific publishing.

\lipsum[1]

% --------------------------------------------
\section{Mathematics}
% --------------------------------------------

\LaTeX{} excels at typesetting mathematical equations.

\subsection{Inline Math}

Einstein's famous equation is $E = mc^2$. The quadratic formula is $x = \frac{-b \pm \sqrt{b^2 - 4ac}}{2a}$.

\subsection{Display Math}

The Gaussian integral:
\begin{equation}
    \int_{-\infty}^{\infty} e^{-x^2} \, dx = \sqrt{\pi}
    \label{eq:gaussian}
\end{equation}

The Euler's identity:
\begin{equation}
    e^{i\pi} + 1 = 0
    \label{eq:euler}
\end{equation}

A system of equations:
\begin{align}
    \nabla \cdot \mathbf{E} &= \frac{\rho}{\varepsilon_0} \\
    \nabla \cdot \mathbf{B} &= 0 \\
    \nabla \times \mathbf{E} &= -\frac{\partial \mathbf{B}}{\partial t} \\
    \nabla \times \mathbf{B} &= \mu_0 \mathbf{J} + \mu_0 \varepsilon_0 \frac{\partial \mathbf{E}}{\partial t}
\end{align}

% --------------------------------------------
\section{Lists}
% --------------------------------------------

\subsection{Itemized List}

\begin{itemize}
    \item First item
    \item Second item
    \item Third item with sub-items:
    \begin{itemize}
        \item Sub-item A
        \item Sub-item B
    \end{itemize}
\end{itemize}

\subsection{Enumerated List}

\begin{enumerate}
    \item Install a \LaTeX{} distribution (MiKTeX, TeX Live, or MacTeX)
    \item Install the LaTeX Workshop extension in VS Code
    \item Create a \texttt{.tex} file
    \item Write your document
    \item Compile and preview
\end{enumerate}

% --------------------------------------------
\section{Tables}
% --------------------------------------------

\begin{table}[h]
    \centering
    \caption{Sample Table}
    \label{tab:sample}
    \begin{tabular}{|l|c|r|}
        \hline
        \textbf{Left} & \textbf{Center} & \textbf{Right} \\
        \hline
        Item 1 & Value 1 & 100 \\
        Item 2 & Value 2 & 200 \\
        Item 3 & Value 3 & 300 \\
        \hline
    \end{tabular}
\end{table}

% --------------------------------------------
\section{Figures}
% --------------------------------------------

To include a figure, use the following template (uncomment when you have an image):

% \begin{figure}[h]
%     \centering
%     \includegraphics[width=0.8\textwidth]{your-image.png}
%     \caption{Your caption here}
%     \label{fig:sample}
% \end{figure}

% --------------------------------------------
\section{Code Listings}
% --------------------------------------------

Here's an example of a Python code listing:

\begin{lstlisting}[language=Python, caption={Hello World in Python}]
def hello_world():
    """A simple hello world function."""
    print("Hello, World!")
    
if __name__ == "__main__":
    hello_world()
\end{lstlisting}

% --------------------------------------------
\section{Cross-References}
% --------------------------------------------

You can reference equations, tables, and figures:
\begin{itemize}
    \item Equation~\ref{eq:gaussian} shows the Gaussian integral
    \item Equation~\ref{eq:euler} is Euler's identity
    \item Table~\ref{tab:sample} demonstrates table formatting
\end{itemize}

% --------------------------------------------
\section{Conclusion}
% --------------------------------------------

This template demonstrates the basic features of \LaTeX{}. Happy typesetting!

\lipsum[2]

% --------------------------------------------
\section{Citations and Bibliography}
% --------------------------------------------

Here are some example citations from the bibliography:

\begin{itemize}
    \item The foundational work on \TeX{} by Knuth~\cite{knuth1984}
    \item The \LaTeX{} manual by Lamport~\cite{lamport1994}
    \item Einstein's famous paper~\cite{einstein1905}
    \item The official \LaTeX{} Project website~\cite{latexproject}
    \item A sample conference paper~\cite{sample2023}
\end{itemize}

% Bibliography
\bibliographystyle{plain}
\bibliography{references}

\end{document}
